\chapter*{ประวัติผู้จัดทำ}
\addcontentsline{toc}{chapter}{ประวัติผู้จัดทำ}

\begin{tcolorbox}[
    enhanced,
    colback=white,
    colframe=rbruNavy,
    arc=3mm,
    boxrule=1pt,
    leftrule=5mm,
    top=12pt, bottom=12pt, left=14pt, right=14pt
]
    \begin{minipage}{0.30\textwidth}
        \centering
        \begin{tikzpicture}
            \draw[circle, fill=rbruNavy!10, draw=rbruGold, line width=2.5pt, drop shadow={opacity=0.3}] (0,0) circle (1.8cm);
            \node at (0,0.3) {\fontsize{14}{16}\selectfont\bfseries\color{rbruNavy} ผศ.ดร.};
            \node at (0,-0.3) {\fontsize{12}{14}\selectfont\bfseries\color{rbruEmerald} ชีวะ ทัศนา};
        \end{tikzpicture}
    \end{minipage}
    \hfill
    \begin{minipage}{0.66\textwidth}
        {\fontsize{16}{20}\selectfont\bfseries\color{rbruNavy} ผู้ช่วยศาสตราจารย์ ดร.ชีวะ ทัศนา\par}
        \vspace{2mm}
        {\fontsize{11}{14}\selectfont\bfseries\color{rbruGold} หน่วยวิจัยปัญญาประดิษฐ์เพื่อเกษตรดิจิทัล\par}
        {\fontsize{10}{13}\selectfont\color{rbruSlate} หลักสูตร คบ.ฟิสิกส์ คณะวิทยาศาสตร์และเทคโนโลยี\\มหาวิทยาลัยราชภัฏรำไพพรรณี\par}
        \vspace{2mm}
        {\fontsize{9.5}{12}\selectfont\color{rbruBlue} \textbf{อีเมล} chewa.t@rbru.ac.th \quad \textbf{สถิติการวิจัย} Google Scholar `jJz3B6AAAAAJ`\par}
    \end{minipage}
\end{tcolorbox}

\vspace{0.6cm}
{\color{rbruGold}\hrule height 1.2pt}
\vspace{0.6cm}

\begin{minipage}[t]{0.48\textwidth}
    {\fontsize{12}{15}\selectfont\bfseries\color{rbruNavy} ประวัติการศึกษา\par}
    \vspace{2mm}
    \begin{itemize}[leftmargin=1.2em, itemsep=4pt]
        \item \textbf{ปร.ด. (ฟิสิกส์ประยุกต์)}\\สถาบันเทคโนโลยีพระจอมเกล้าเจ้าคุณทหารลาดกระบัง (KMITL)
        \item \textbf{วท.ม. (ฟิสิกส์ประยุกต์)}\\สถาบันเทคโนโลยีพระจอมเกล้าเจ้าคุณทหารลาดกระบัง (KMITL)
        \item \textbf{วท.บ. ฟิสิกส์ (คอมพิวเตอร์และอิเล็กทรอนิกส์)}\\มหาวิทยาลัยนเรศวร
    \end{itemize}

    \vspace{4mm}
    {\fontsize{12}{15}\selectfont\bfseries\color{rbruNavy} ภาระงานสอนและการพัฒนานักศึกษา\par}
    \vspace{2mm}
    \begin{itemize}[leftmargin=1.2em, itemsep=3pt]
        \item ฟิสิกส์เกษตรดิจิทัลและสมองกลฝังตัว
        \item ปัญญาประดิษฐ์ประมวลผลบนขอบ (Edge AI)
        \item อิเล็กทรอนิกส์และการประมวลผลสัญญาณ
    \end{itemize}
\end{minipage}
\hfill
\begin{minipage}[t]{0.48\textwidth}
    {\fontsize{12}{15}\selectfont\bfseries\color{rbruNavy} ความเชี่ยวชาญและงานวิจัยที่สนใจ\par}
    \vspace{2mm}
    \begin{itemize}[leftmargin=1.2em, itemsep=4pt]
        \item \textbf{ฟิสิกส์เกษตรดิจิทัล (Digital Agriphysics)} การบูรณาการทฤษฎีทางฟิสิกส์เคมีเข้ากับการตรวจวัดสุขภาพดิน
        \item \textbf{ปัญญาประดิษฐ์บนขอบ (Embedded Edge AI / TinyML)} โครงข่ายประสาทเทียมและอัลกอริทึม OLS ประมวลผลบนชิปสมาร์ทโฟน
        \item \textbf{ระบบเซนเซอร์อัจฉริยะ (Smart Sensor Networks)} การสื่อสาร RS485 Modbus RTU และการเชื่อมต่อผ่านพอร์ต USB-C
        \item \textbf{การเกษตรแม่นยำสูง (High-Precision Agriculture)} การจัดการดินและธาตุอาหารในสวนทุเรียนแปลงใหญ่ภาคตะวันออก
    \end{itemize}
\end{minipage}

\vspace{1.0cm}
\begin{center}
    {\fontsize{9}{12}\selectfont\itshape\color{rbruSlate} เอกสารนี้จัดทำขึ้นโดย ผศ.ดร.ชีวะ ทัศนา หน่วยวิจัยปัญญาประดิษฐ์เพื่อเกษตรดิจิทัล คณะวิทยาศาสตร์และเทคโนโลยี มหาวิทยาลัยราชภัฏรำไพพรรณี\par}
\end{center}

\clearpage
